%%=============================================================================
%% Vragenlijst survey
%%=============================================================================

\chapter{Vragenlijst survey}%
\label{bijl:surveyanalyse}

Deze bijlage geeft de volledige vragenlijst weer van de kwantitatieve survey
beschreven in Hoofdstuk~\ref{ch:surveyanalyse}. De vragenlijst werd opgesteld
in Google Forms en is opgedeeld in vijf blokken: profielvragen, vragen over
huidige planningstools, designkeuzes via A/B-vergelijkingen,
prototype-evaluatie op basis van Likert-schalen en een open eindvraag.

\section*{Blok 1: Profielvragen}

\begin{enumerate}
  \item Heb je een diagnose of vermoeden van ADHD?
  \begin{itemize}
    \item Ja, ik heb een officiële diagnose
    \item Ik vermoed het sterk (geen officiële diagnose)
    \item Nee, maar ik herken mij in planningsproblemen
  \end{itemize}
  \item Hoe zou je jouw relatie met deadlines omschrijven? \textit{(meerdere antwoorden mogelijk)}
  \begin{itemize}
    \item Ik begin ruim op tijd en werk gespreid
    \item Ik begin pas als de deadline dichtbij is
    \item Ik weet dat ik moet beginnen maar geraak er niet aan toe
    \item Ik mis deadlines regelmatig, ook al wil ik dat niet
  \end{itemize}
  \item Wat is het moeilijkste moment in een planningsproces voor jou?
  \begin{itemize}
    \item Inschatten hoeveel tijd iets kost
    \item Aan iets beginnen, ook als ik wil
    \item Verdergaan na een onderbreking of gemiste deadline
  \end{itemize}
  \item Hoe reageer je als je een deadline hebt gemist?
  \begin{itemize}
    \item Ik herplan meteen en ga verder
    \item Ik doe het op het allerlaatste moment alsnog
    \item Ik stel herplannen uit door schaamte of stress
    \item Ik geef de taak helemaal op
  \end{itemize}
\end{enumerate}

\section*{Blok 2: Vragen over huidige planningstools}

\begin{enumerate}
  \setcounter{enumi}{4}
  \item Welke planningstools heb je al geprobeerd? \textit{(meerdere antwoorden mogelijk)}
  \begin{itemize}
    \item Google Calendar
    \item Microsoft To Do
    \item Todoist
    \item Notion
    \item Tiimo
    \item Goblin.tools
    \item Papieren agenda of notitieboek
    \item Ik gebruik geen planningssysteem
  \end{itemize}
  \item Hoe lang houd je een planningsapp gemiddeld vol?
  \begin{itemize}
    \item Minder dan een week
    \item Één tot vier weken
    \item Één tot drie maanden
    \item Langer dan drie maanden
    \item Ik gebruik nooit planningsapps
  \end{itemize}
  \item Wat is de voornaamste reden dat je met een planningsapp stopt?
  \begin{itemize}
    \item Te veel moeite om alles in te voeren
    \item De app helpt niet bij het daadwerkelijk starten
    \item Notificaties worden genegeerd of storen
    \item Na een gemiste dag of deadline stop ik gewoon
  \end{itemize}
  \item In hoeverre vind je dat bestaande planningsapps rekening houden met
  hoe jouw brein werkt? \textit{(Likertschaal 1 = helemaal niet, 5 = volledig)}
\end{enumerate}

\section*{Blok 3: Designkeuzes via A/B-vergelijkingen}

\begin{enumerate}
  \setcounter{enumi}{8}
  \item Hoe wil je een taak invoeren?
  \begin{itemize}
    \item Gestructureerd formulier: titel, deadline en prioriteit tegelijk
    invullen. Voorbeeld: velden voor naam, datum en grootte meteen zichtbaar.
    \item Vrij typen: je typt wat je wil doen, hoe vaag ook. De app vraagt
    daarna pas om details. Voorbeeld: ``verslag schrijven'' $\rightarrow$ app
    vraagt: deadline? hoe groot?
  \end{itemize}
  \item Hoe moet de AI een grote taak opsplitsen?
  \begin{itemize}
    \item Logische stappen: de AI maakt stappen op basis van wat logisch
    samenhangt (bv. ``lezen'', ``samenvatten'', ``schrijven'').
    \item Vaste tijdsblokken: de AI verdeelt de taak in blokken van vaste
    duur (bv. 25 min, 25 min, 30 min).
  \end{itemize}
  \item Wat moet er gebeuren als je een deadline mist?
  \begin{itemize}
    \item Volledig overzicht: de app toont alles wat gemist werd met de
    optie om alles in één keer te herplannen.
    \item Neutrale herstart: de app toont de volgende ochtend: ``Één ding
    bleef liggen, wil je het vandaag inplannen?''
  \end{itemize}
  \item Hoe wil je herinnerd worden aan een taak?
  \begin{itemize}
    \item Widget op startscherm: een permanent zichtbaar overzichtje op je
    startscherm, zonder actieve melding.
    \item Dagelijkse melding: één pushmelding per dag met een overzicht van
    wat er die dag staat.
  \end{itemize}
  \item Hoe moet de dagweergave werken?
  \begin{itemize}
    \item Tijdlijn: taken op een visuele tijdslijn op basis van uur (bv.
    10u--11u: lezen, 13u: vergadering).
    \item Prioriteitenlijst: taken gesorteerd op prioriteit, zonder vaste
    tijden. Jij bepaalt wanneer je wat doet.
  \end{itemize}
  \item Hoe moet de app reageren als je zegt dat je vastloopt bij een taak?
  \begin{itemize}
    \item Verder opsplitsen: de app breekt de huidige stap op in nog
    kleinere micro-stappen.
    \item Alternatieve taak voorstellen: de app stelt een andere taak voor
    die je op dat moment wél kan doen.
  \end{itemize}
  \item Wat mag een app absoluut NIET doen? \textit{(selecteer alles dat van toepassing is)}
  \begin{itemize}
    \item Mij zeggen wat ik nu op dit moment moet doen
    \item Mij herinneren aan wat ik niet gedaan heb
    \item Veel invoer vragen voor ik iets kan doen
    \item Mij straffen (visueel of via tellers) voor gemiste taken
    \item Niets van bovenstaande stoort mij
  \end{itemize}
\end{enumerate}

\section*{Blok 4: Prototype-evaluatie (Likertschaal 1--5, helemaal niet
akkoord -- volledig akkoord)}

\begin{enumerate}
  \setcounter{enumi}{15}
  \item Ik vond het prototype makkelijk te gebruiken.
  \item De interface voelde rustig en overzichtelijk aan.
  \item De manier waarop de app reageert op een gemiste deadline voelt
  gepast aan voor mij.
  \item Ik zou deze app willen gebruiken als hij beschikbaar was op mijn
  telefoon.
\end{enumerate}

\section*{Blok 5: Open eindvragen}

\begin{enumerate}
  \setcounter{enumi}{19}
  \item Welke functie vond je het meest waardevol in het prototype?
  \begin{itemize}
    \item De dagweergave met kleurgecodeerde taken
    \item Het opsplitsen van taken in stappen (taakdetailscherm)
    \item De patroonherkenning en inzichten
    \item Het scherm bij een gemiste deadline --- geen straf, wel hulp
    \item Geen van bovenstaande sprak me aan
  \end{itemize}
  \item Wat zou jou het meest tegenhouden om deze app dagelijks te
  gebruiken?
  \begin{itemize}
    \item Ik mis een echte mobiele app (geen browser)
    \item Ik zou snel de gewoonte verliezen om hem te openen
    \item Het invoeren van taken kost nog te veel moeite
    \item De navigatie is niet duidelijk genoeg (te weinig tabbladen)
    \item De stappen die de AI genereert kloppen niet altijd
    \item Niets, ik zou hem graag gebruiken
  \end{itemize}
  \item Is er iets dat jij absoluut in een planningsapp zou willen dat hier
  niet in zit, of iets dat je anders zou aanpakken? \textit{(open vraag)}
\end{enumerate}
